\documentclass[11pt,a4paper]{article}
\usepackage{examgrid}
\setdoc{System Design and ERD Specification}{EG-DOC-02}

\begin{document}
\nmamitcover{System Design and\\ER Model Specification}{Document 02 of 05 \,\textbullet\, High-Level and Low-Level Design}{\small Department of Computer Science and Engineering\\[6pt]November 2026}

\pagenumbering{roman}
\doccontrol{EG-DOC-02}{Approved baseline}
\newpage
\tableofcontents
\listoffigures
\listoftables
\newpage
\pagenumbering{arabic}

\section{Introduction}
\subsection{Purpose}
This document describes how ExamGrid is built: the overall architecture, the Spring MVC layers, the package and class structure, the MySQL database with its entity-relationship model, the main runtime interactions and the user interface navigation. It refines the requirements of Document 01 (SRS) into a design that the code in the repository follows one-to-one.

\subsection{Design Goals}
\begin{enumerate}
  \item \textbf{Separation of concerns} -- presentation, business rules and persistence live in separate layers.
  \item \textbf{Single source of truth} -- seating and duty rules exist only in two service classes.
  \item \textbf{Data integrity in the database} -- uniqueness and foreign keys are enforced by MySQL, not only by Java code.
  \item \textbf{Simplicity} -- a single deployable Spring Boot JAR, readable code suitable for a B.Tech viva.
\end{enumerate}

\section{Architectural Overview}
ExamGrid uses a classic \textbf{layered Model--View--Controller (MVC)} architecture on top of Spring Boot. Each HTTP request first passes through the Spring Security filter chain, then reaches a controller. Controllers only translate HTTP to Java method calls and choose a view or a JSON response; all decisions are taken in the service layer; services use Spring Data JPA repositories to read and write MySQL.

\begin{figure}[H]\centering
\resizebox{0.85\linewidth}{!}{\begin{tikzpicture}[
  font=\small,
  layer/.style={draw=egnavy,fill=eglight,rounded corners=3pt,minimum width=9.6cm,minimum height=1.05cm,align=center},
  comp/.style={draw=egnavy,fill=white,rounded corners=2pt,minimum height=0.7cm,align=center,font=\footnotesize},
  side/.style={draw=eggold,fill=eggold!15,rounded corners=3pt,minimum width=3.0cm,minimum height=1.05cm,align=center,font=\footnotesize},
  arr/.style={-{Stealth[length=2.5mm]},thick,egnavy}]

\node[layer,fill=white] (client) {\textbf{Client Tier} -- Web browser (Admin / Faculty / Student), Postman, curl};
\node[layer,below=0.55cm of client,fill=eggold!20,draw=eggold] (sec) {\textbf{Spring Security Filter Chains}\\[-1pt]\footnotesize API chain (/api/**, HTTP Basic) \quad Web chain (form login, session)};
\node[layer,below=0.55cm of sec,minimum height=1.6cm] (ctl) {};
\node[anchor=north] at (ctl.north) {\textbf{Controller Layer}};
\node[comp,anchor=south west] at ([xshift=0.25cm,yshift=0.15cm]ctl.south west) {Web Controllers\\(@Controller)};
\node[comp,anchor=south] at ([yshift=0.15cm]ctl.south) {REST Controllers\\(@RestController)};
\node[comp,anchor=south east] at ([xshift=-0.25cm,yshift=0.15cm]ctl.south east) {@ControllerAdvice\\handlers};
\node[layer,below=0.55cm of ctl,minimum height=1.6cm] (svc) {};
\node[anchor=north] at (svc.north) {\textbf{Service Layer} (@Service, @Transactional)};
\node[comp,anchor=south west] at ([xshift=0.2cm,yshift=0.15cm]svc.south west) {SeatingAllocation\\Service};
\node[comp,anchor=south] at ([xshift=-0.15cm,yshift=0.15cm]svc.south) {InvigilationAllocation\\Service};
\node[comp,anchor=south east] at ([xshift=-0.2cm,yshift=0.15cm]svc.south east) {CRUD + ReportPdf\\Services};
\node[layer,below=0.55cm of svc] (repo) {\textbf{Repository Layer} -- Spring Data JPA interfaces (JpaRepository)};
\node[layer,below=0.55cm of repo] (orm) {\textbf{Persistence} -- Hibernate ORM \,+\, JDBC (MySQL Connector/J)};
\node[layer,below=0.55cm of orm,fill=egnavy,text=white] (db) {\textbf{MySQL 8} -- examgrid\_db};

\node[side,right=0.6cm of ctl] (view) {\textbf{View}\\Thymeleaf templates\\Bootstrap 5, JS Fetch};
\node[side,right=0.6cm of svc] (pdf) {\textbf{OpenPDF}\\Seating charts\\Duty rosters};
\node[side,right=0.6cm of repo] (ent) {\textbf{Model}\\JPA Entities\\DTO records};

\draw[arr] (client) -- node[right,font=\scriptsize]{HTTP request} (sec);
\draw[arr] (sec) -- (ctl);
\draw[arr] (ctl) -- (svc);
\draw[arr] (svc) -- (repo);
\draw[arr] (repo) -- (orm);
\draw[arr] (orm) -- (db);
\draw[arr,dashed] (ctl) -- (view);
\draw[arr,dashed] (svc) -- (pdf);
\draw[arr,dashed] (repo) -- (ent);
\draw[arr,eggold] (view.north) |- node[pos=0.25,right,font=\scriptsize]{HTML / JSON} (client.east);
\end{tikzpicture}
}
\caption{Layered MVC architecture of ExamGrid}\label{fig:arch}
\end{figure}

\subsection{Layer Responsibilities}
\begin{tabularx}{\linewidth}{|l|L{3.3cm}|X|}
\hline\headrow\thead{Layer} & \thead{Spring stereotype} & \thead{Responsibility}\\\hline
View & Thymeleaf templates & Render HTML from the model; Bootstrap layout; JavaScript Fetch calls for seat lookup and roster.\\\hline
Controller (web) & \texttt{@Controller} & Handle form posts, run Bean Validation (\texttt{@Valid}), add model attributes, redirect with flash messages.\\\hline
Controller (API) & \texttt{@RestController} & Expose JSON endpoints under \texttt{/api}; return proper HTTP status codes.\\\hline
Error handling & \texttt{@ControllerAdvice}, \texttt{@RestControllerAdvice} & Convert exceptions into an error page (web) or an \texttt{ErrorResponse} JSON (API).\\\hline
Service & \texttt{@Service}, \texttt{@Transactional} & Business rules, allocation engines, PDF creation, account creation.\\\hline
Repository & \texttt{JpaRepository} interfaces & Derived queries and JPQL with \texttt{join fetch}; bulk deletes.\\\hline
Model & \texttt{@Entity}, Java records & JPA entities mapped to tables; immutable DTO records for API responses.\\\hline
Security & \texttt{SecurityFilterChain} beans & Authentication (form login, HTTP Basic), role-based authorization, sessions, CSRF.\\\hline
\end{tabularx}

\subsection{Why Two Security Filter Chains}
Browser users need a login form, a session cookie and CSRF protection. API clients such as Postman and curl cannot easily handle a login form, so the \texttt{/api/**} chain uses HTTP Basic authentication and disables CSRF. The chains are ordered: \texttt{@Order(1)} for the API matcher and \texttt{@Order(2)} for everything else.

\section{Package Structure}
\begin{lstlisting}[style=plain]
com.examgrid
 |-- ExamGridApplication.java          (Spring Boot entry point)
 |-- config/       DataSeeder          (sample data on first start)
 |-- controller/
 |    |-- api/     StudentApiController, FacultyApiController, HallApiController,
 |    |            ExamApiController, AllocationApiController,
 |    |            SeatLookupApiController, DutyApiController
 |    `-- web/     HomeController, AdminDashboardController, StudentWebController,
 |                 FacultyWebController, HallWebController, ExamWebController,
 |                 AllocationWebController, ReportController, PortalController
 |-- dto/          SeatInfo, DutyInfo, SeatingResult, InvigilationResult,
 |                 ExamSlot, HallLayout, ErrorResponse      (Java records)
 |-- entity/       Student, Faculty, Hall, Exam, SeatAllocation,
 |                 InvigilationDuty, AppUser, Role
 |-- exception/    ResourceNotFoundException, DuplicateResourceException,
 |                 InsufficientCapacityException, AllocationException,
 |                 InvalidRequestException, ApiExceptionHandler, WebExceptionHandler
 |-- repository/   7 Spring Data JPA repositories
 |-- security/     SecurityConfig, AppUserDetailsService
 `-- service/      SeatingAllocationService, InvigilationAllocationService,
                   StudentService, FacultyService, HallService, ExamService,
                   UserAccountService, ReportPdfService
\end{lstlisting}

\section{Class Design}
\subsection{Domain Classes}
\begin{figure}[H]\centering
\begin{tikzpicture}[font=\scriptsize,
  cls/.style={rectangle split,rectangle split parts=3,draw=egnavy,thick,rectangle split part fill={eglight,white,white},text width=4.1cm,align=left,inner sep=3pt},
  as/.style={thick,egnavy}]
\node[cls] (st) at (0,0) {\centering\textbf{Student}\nodepart{two}-id: Long\\-usn: String\\-name: String\\-branch: String\\-semester: int\\-email: String\nodepart{three}+getters / setters};
\node[cls] (ex) at (0,-4.6) {\centering\textbf{Exam}\nodepart{two}-id: Long\\-subjectCode, subjectName\\-branch: String; semester: int\\-examDate: LocalDate\\-startTime, endTime: LocalTime\nodepart{three}+getters / setters};
\node[cls] (sa) at (5.4,-2.3) {\centering\textbf{SeatAllocation}\nodepart{two}-id: Long\\-exam: Exam\\-student: Student\\-hall: Hall\\-examDate, startTime\\-rowNumber, columnNumber: int\\-seatNumber: String\nodepart{three}+SeatAllocation(exam, student, hall, row, col)};
\node[cls] (ha) at (10.8,-0.4) {\centering\textbf{Hall}\nodepart{two}-id: Long\\-hallName, block: String\\-numberOfRows: int\\-numberOfColumns: int\nodepart{three}+getCapacity(): int};
\node[cls] (du) at (10.8,-4.9) {\centering\textbf{InvigilationDuty}\nodepart{two}-id: Long\\-faculty: Faculty\\-hall: Hall\\-examDate: LocalDate\\-startTime, endTime\nodepart{three}+overlaps(date, from, to): boolean};
\node[cls] (fa) at (5.4,-7.0) {\centering\textbf{Faculty}\nodepart{two}-id: Long\\-facultyCode, name\\-department, email\\-maxDuties: int = 6\nodepart{three}+getters / setters};
\node[cls] (us) at (0,-8.6) {\centering\textbf{AppUser}\nodepart{two}-id: Long\\-username, password\\-role: Role\\-enabled: boolean\nodepart{three}+getters / setters};
\draw[as] (st.east) -- node[above,pos=0.15]{1} node[above,pos=0.85]{*} (sa.north west);
\draw[as] (ex.east) -- node[below,pos=0.15]{1} node[below,pos=0.85]{*} (sa.south west);
\draw[as] (ha.west) -- node[above,pos=0.15]{1} node[above,pos=0.85]{*} (sa.east);
\draw[as] (ha.south) -- node[right,pos=0.15]{1} node[right,pos=0.85]{*} (du.north);
\draw[as] (fa.east) -- node[below,pos=0.15]{1} node[below,pos=0.85]{*} (du.west);
\end{tikzpicture}
\caption{Class diagram of JPA entities}\label{fig:class}
\end{figure}

\subsection{Service Classes}
\begin{tabularx}{\linewidth}{|L{4.4cm}|X|}
\hline\headrow\thead{Service} & \thead{Public operations}\\\hline
SeatingAllocationService & \texttt{allocateSeats(date, startTime)}, \texttt{findSeatsByUsn(usn)}, \texttt{findSlotSeats(date, time)}, \texttt{buildHallLayouts(date, time)}\\\hline
InvigilationAllocationService & \texttt{allocateInvigilators(date, startTime)}, \texttt{findAllDuties()}, \texttt{findDutiesByDate(date)}, \texttt{findDutiesOfFaculty(code)}, \texttt{workloadSummary()}\\\hline
StudentService / FacultyService & \texttt{findAll}, \texttt{findById}, \texttt{create}, \texttt{update}, \texttt{delete}; create also makes a login account\\\hline
HallService & CRUD plus \texttt{totalCapacity()}\\\hline
ExamService & CRUD, \texttt{findTimingProblem(exam)}, \texttt{findSlots()}\\\hline
UserAccountService & \texttt{createDefaultAccount}, \texttt{deleteAccount}; BCrypt encoding\\\hline
ReportPdfService & \texttt{seatingChart(date, time)}, \texttt{dutyRoster(date)} returning \texttt{byte[]}\\\hline
\end{tabularx}

All services receive their dependencies through constructor injection, which makes them easy to test and avoids field injection.

\section{Database Design}
\subsection{Entity-Relationship Diagram}
\begin{figure}[H]\centering
\resizebox{\linewidth}{!}{\begin{tikzpicture}[
  font=\scriptsize,
  ent/.style={rectangle split,rectangle split parts=2,draw=egnavy,thick,rounded corners=2pt,
              rectangle split part fill={egnavy,white},text width=3.25cm,align=left,inner sep=3pt},
  rel/.style={diamond,draw=eggold,fill=eggold!20,thick,aspect=2.2,inner sep=1pt,font=\scriptsize\itshape},
  ln/.style={thick,egnavy!80}]

\node[ent] (student) at (0,0) {\textcolor{white}{\textbf{STUDENTS}}\nodepart{two}
  \underline{id} BIGINT PK\\ usn VARCHAR(15) UK\\ name VARCHAR(80)\\ branch VARCHAR(10)\\ semester INT\\ email VARCHAR(100)};
\node[ent] (exam) at (0,-6.2) {\textcolor{white}{\textbf{EXAMS}}\nodepart{two}
  \underline{id} BIGINT PK\\ subject\_code VARCHAR(15)\\ subject\_name VARCHAR(100)\\ branch VARCHAR(10)\\ semester INT\\ exam\_date DATE\\ start\_time TIME\\ end\_time TIME};
\node[ent] (seat) at (6,-3.1) {\textcolor{white}{\textbf{SEAT\_ALLOCATIONS}}\nodepart{two}
  \underline{id} BIGINT PK\\ exam\_id FK\\ student\_id FK\\ hall\_id FK\\ exam\_date DATE\\ start\_time TIME\\ seat\_row INT\\ seat\_column INT\\ seat\_number VARCHAR(20)};
\node[ent] (hall) at (12,-3.1) {\textcolor{white}{\textbf{HALLS}}\nodepart{two}
  \underline{id} BIGINT PK\\ hall\_name VARCHAR(30) UK\\ block VARCHAR(30)\\ total\_rows INT\\ total\_columns INT};
\node[ent] (duty) at (12,-8.4) {\textcolor{white}{\textbf{INVIGILATION\_DUTIES}}\nodepart{two}
  \underline{id} BIGINT PK\\ faculty\_id FK\\ hall\_id FK\\ exam\_date DATE\\ start\_time TIME\\ end\_time TIME};
\node[ent] (faculty) at (5.5,-9.6) {\textcolor{white}{\textbf{FACULTY}}\nodepart{two}
  \underline{id} BIGINT PK\\ faculty\_code VARCHAR(15) UK\\ name VARCHAR(80)\\ department VARCHAR(10)\\ email VARCHAR(100)\\ max\_duties INT};
\node[ent] (user) at (12,1.2) {\textcolor{white}{\textbf{APP\_USERS}}\nodepart{two}
  \underline{id} BIGINT PK\\ username VARCHAR(30) UK\\ password VARCHAR(255)\\ role ENUM\\ enabled BIT};

\node[rel] (r1) at (3.4,-0.6) {allotted};
\node[rel] (r2) at (3.4,-6.0) {seats};
\node[rel] (r3) at (9.2,-2.0) {in};
\node[rel] (r4) at (12,-5.9) {covers};
\node[rel] (r5) at (9.0,-9.3) {performs};

\draw[ln] (student.east) -- node[above,pos=0.2]{1} (r1) -- node[above,pos=0.8]{N} (seat.north west);
\draw[ln] (exam.east) -- node[above,pos=0.2]{1} (r2) -- node[above,pos=0.8]{N} (seat.south west);
\draw[ln] (hall.west) -- node[above,pos=0.2]{1} (r3) -- node[above,pos=0.8]{N} (seat.east);
\draw[ln] (hall.south) -- node[right,pos=0.2]{1} (r4) -- node[right,pos=0.8]{N} (duty.north);
\draw[ln] (faculty.east) -- node[above,pos=0.2]{1} (r5) -- node[above,pos=0.8]{N} (duty.west);
\draw[ln,dashed,eggrey] (user.west) -- node[above,font=\tiny,align=center]{username = usn}  (student.north east);
\node[font=\tiny,eggrey,align=center,below=1pt of user] {login account: username = USN (student)\\or faculty\_code (faculty) or ADMIN};
\end{tikzpicture}
}
\caption{ER diagram of \texttt{examgrid\_db}}\label{fig:er}
\end{figure}

\subsection{Relationships}
\begin{tabularx}{\linewidth}{|l|l|X|}
\hline\headrow\thead{Relationship} & \thead{Cardinality} & \thead{Meaning}\\\hline
Student -- SeatAllocation & 1 : N & A student gets one seat per exam he or she writes.\\\hline
Exam -- SeatAllocation & 1 : N & An exam has one seat row per student writing it.\\\hline
Hall -- SeatAllocation & 1 : N & A hall holds many seats in many slots.\\\hline
Faculty -- InvigilationDuty & 1 : N & A faculty member has many duties over the exam period.\\\hline
Hall -- InvigilationDuty & 1 : N & A hall needs one or more invigilators per slot.\\\hline
AppUser -- Student / Faculty & 1 : 1 (logical) & Linked by \texttt{username = usn} or \texttt{username = faculty\_code}; kept separate so that credentials are not mixed with academic data.\\\hline
\end{tabularx}

SeatAllocation is an \emph{associative entity} that resolves the many-to-many relationship between Student, Exam and Hall and adds the seat position attributes.

\subsection{Data Dictionary}
\newcommand{\ddtable}[2]{%
\begin{table}[H]\centering\small
\caption{Table \texttt{#1}}
\begin{tabularx}{\linewidth}{|l|l|l|X|}
\hline\headrow\thead{Column} & \thead{Type} & \thead{Constraint} & \thead{Description}\\\hline
#2
\end{tabularx}\end{table}}

\ddtable{students}{
id & BIGINT & PK, auto increment & Surrogate key\\\hline
usn & VARCHAR(15) & NOT NULL, UNIQUE & University Seat Number, stored upper case\\\hline
name & VARCHAR(80) & NOT NULL & Full name\\\hline
branch & VARCHAR(10) & NOT NULL & CSE, ISE, ECE, EEE, ME, CV, AIML\\\hline
semester & INT & 1--8 (Bean Validation) & Current semester\\\hline
email & VARCHAR(100) & valid e-mail & Institute e-mail\\\hline}

\ddtable{faculty}{
id & BIGINT & PK & Surrogate key\\\hline
faculty\_code & VARCHAR(15) & NOT NULL, UNIQUE & Staff code, also the login name\\\hline
name & VARCHAR(80) & NOT NULL & Full name with title\\\hline
department & VARCHAR(10) & NOT NULL & Department code\\\hline
email & VARCHAR(100) & valid e-mail & Institute e-mail\\\hline
max\_duties & INT & 1--20, default 6 & Upper limit of invigilation duties\\\hline}

\ddtable{halls}{
id & BIGINT & PK & Surrogate key\\\hline
hall\_name & VARCHAR(30) & NOT NULL, UNIQUE & e.g. LH-101\\\hline
block & VARCHAR(30) & NOT NULL & Building / block name\\\hline
total\_rows & INT & 1--30 & Rows of benches from the blackboard backwards\\\hline
total\_columns & INT & 1--20 & Seats in one row\\\hline}

\ddtable{exams}{
id & BIGINT & PK & Surrogate key\\\hline
subject\_code & VARCHAR(15) & NOT NULL & e.g. CS3603\\\hline
subject\_name & VARCHAR(100) & NOT NULL & Subject title\\\hline
branch, semester & VARCHAR(10), INT & NOT NULL & Class that writes the exam\\\hline
exam\_date & DATE & NOT NULL & Exam day\\\hline
start\_time, end\_time & TIME & NOT NULL, end $>$ start & Exam timing; exams with same date and start form a slot\\\hline}

\ddtable{seat\_allocations}{
id & BIGINT & PK & Surrogate key\\\hline
exam\_id & BIGINT & FK $\rightarrow$ exams & Exam being written\\\hline
student\_id & BIGINT & FK $\rightarrow$ students & Seated student\\\hline
hall\_id & BIGINT & FK $\rightarrow$ halls & Allotted hall\\\hline
exam\_date, start\_time & DATE, TIME & NOT NULL & Copied from exam to identify the slot quickly\\\hline
seat\_row, seat\_column & INT & NOT NULL & 1-based position (row 1 is nearest the blackboard)\\\hline
seat\_number & VARCHAR(20) & NOT NULL & Printable label, e.g. SH-01-C4\\\hline
\multicolumn{4}{|l|}{UNIQUE (exam\_id, student\_id); UNIQUE (hall\_id, exam\_date, start\_time, seat\_row, seat\_column)}\\\hline}

\ddtable{invigilation\_duties}{
id & BIGINT & PK & Surrogate key\\\hline
faculty\_id & BIGINT & FK $\rightarrow$ faculty & Invigilator\\\hline
hall\_id & BIGINT & FK $\rightarrow$ halls & Hall supervised\\\hline
exam\_date & DATE & NOT NULL & Duty date\\\hline
start\_time, end\_time & TIME & NOT NULL & Duty window\\\hline}

\ddtable{app\_users}{
id & BIGINT & PK & Surrogate key\\\hline
username & VARCHAR(30) & NOT NULL, UNIQUE & ADMIN, USN or faculty code\\\hline
password & VARCHAR(255) & NOT NULL & BCrypt hash (60 chars)\\\hline
role & VARCHAR(10) & ENUM as string & ADMIN, FACULTY or STUDENT\\\hline
enabled & BIT & default 1 & Account active flag\\\hline}

\subsection{Normalization}
\begin{itemize}
  \item \textbf{1NF}: every column holds a single atomic value; there are no repeating groups (seat position is two integer columns, not a list).
  \item \textbf{2NF}: every table has a single-column surrogate key, so no partial dependency is possible.
  \item \textbf{3NF}: non-key columns depend only on the key. Hall capacity is \emph{not stored} because it is derived from rows $\times$ columns.
  \item \textbf{Controlled denormalization}: \texttt{exam\_date} and \texttt{start\_time} are copied into \texttt{seat\_allocations} so that the unique constraint ``one student per physical seat per slot'' can be enforced directly by MySQL and slot queries need no join. The copy is written only by the allocation engine, so it cannot drift.
\end{itemize}

\subsection{Indexes}
MySQL automatically creates indexes for primary keys, unique constraints and foreign keys. The two most frequent queries -- seat lookup by USN and loading one slot -- use the \texttt{students.usn} unique index and the composite unique index on \texttt{(hall\_id, exam\_date, start\_time, \ldots)} respectively.

\subsection{DDL}
The complete script is in \texttt{docs/schema.sql}. Hibernate generates the same tables automatically with \texttt{spring.jpa.hibernate.ddl-auto=update}.
\begin{lstlisting}[style=sql]
CREATE TABLE seat_allocations (
    id           BIGINT AUTO_INCREMENT PRIMARY KEY,
    exam_id      BIGINT      NOT NULL,
    student_id   BIGINT      NOT NULL,
    hall_id      BIGINT      NOT NULL,
    exam_date    DATE        NOT NULL,
    start_time   TIME        NOT NULL,
    seat_row     INT         NOT NULL,
    seat_column  INT         NOT NULL,
    seat_number  VARCHAR(20) NOT NULL,
    CONSTRAINT uk_exam_student UNIQUE (exam_id, student_id),
    CONSTRAINT uk_hall_seat UNIQUE (hall_id, exam_date, start_time, seat_row, seat_column),
    CONSTRAINT fk_seat_exam    FOREIGN KEY (exam_id)    REFERENCES exams (id),
    CONSTRAINT fk_seat_student FOREIGN KEY (student_id) REFERENCES students (id),
    CONSTRAINT fk_seat_hall    FOREIGN KEY (hall_id)    REFERENCES halls (id)
);
\end{lstlisting}
\begin{notebox}[Design decision]
The column names \texttt{seat\_row} and \texttt{total\_rows} are used instead of \texttt{row\_number} and \texttt{rows} because \texttt{ROW\_NUMBER} and \texttt{ROWS} are reserved words in MySQL 8.
\end{notebox}

\subsection{JPA Mapping Choices}
\begin{itemize}
  \item All \texttt{@ManyToOne} associations use \texttt{FetchType.LAZY}; queries that need the related data use \texttt{join fetch} to avoid the N+1 select problem.
  \item \texttt{spring.jpa.open-in-view=false} so that no lazy loading happens in Thymeleaf; services return fully loaded objects or DTO records.
  \item Enums are stored with \texttt{EnumType.STRING} so that reordering the enum never corrupts data.
  \item Java 17 \texttt{LocalDate} and \texttt{LocalTime} map to MySQL \texttt{DATE} and \texttt{TIME}.
\end{itemize}

\section{Dynamic Behaviour}
\subsection{Seat Lookup by USN}
\begin{figure}[H]\centering
\resizebox{\linewidth}{!}{\begin{tikzpicture}[font=\scriptsize,
  head/.style={draw=egnavy,fill=eglight,rounded corners=2pt,minimum width=2.1cm,minimum height=0.65cm,align=center,font=\scriptsize\bfseries},
  msg/.style={-{Stealth[length=2mm]},thick},
  ret/.style={-{Stealth[length=2mm]},dashed}]
\foreach \x/\n/\t in {0/b/Browser\\(seat-lookup.js),3.1/f/Security\\Filter,6.2/c/SeatLookup\\ApiController,9.3/s/SeatingAllocation\\Service,12.4/r/SeatAllocation\\Repository,15.2/d/MySQL}{
  \node[head] (\n) at (\x,0) {\t};
  \draw[eggrey,dashed] (\n.south) -- ++(0,-7.1);
}
\draw[msg] (0,-1.0) -- node[above]{GET /api/seats/NNM24CS224} (3.1,-1.0);
\draw[msg] (3.1,-1.5) -- node[above]{permitAll} (6.2,-1.5);
\draw[msg] (6.2,-2.0) -- node[above]{findSeatsByUsn(usn)} (9.3,-2.0);
\draw[msg] (9.3,-2.5) -- node[above]{existsByUsnIgnoreCase} (12.4,-2.5);
\draw[msg] (12.4,-2.9) -- node[above]{SELECT} (15.2,-2.9);
\draw[msg] (9.3,-3.5) -- node[above]{findByUsn (join fetch)} (12.4,-3.5);
\draw[msg] (12.4,-3.9) -- node[above]{SELECT ... JOIN} (15.2,-3.9);
\draw[ret] (15.2,-4.4) -- node[above]{rows} (12.4,-4.4);
\draw[ret] (12.4,-4.8) -- node[above]{List$<$SeatAllocation$>$} (9.3,-4.8);
\draw[ret] (9.3,-5.3) -- node[above]{List$<$SeatInfo$>$ (DTO)} (6.2,-5.3);
\draw[ret] (6.2,-5.8) -- node[above]{200 OK, JSON array} (0,-5.8);
\draw[msg] (0,-6.3) -- ++(0.6,0) -- ++(0,-0.35) node[right]{render seat cards (DOM)} -- ++(-0.6,0);
\node[draw=red!60!black,fill=red!6,font=\tiny,align=left,text width=3.1cm] at (8.0,-6.75) {alt: unknown USN $\rightarrow$ ResourceNotFound\-Exception $\rightarrow$ ApiExceptionHandler $\rightarrow$ 404 JSON};
\end{tikzpicture}
}
\caption{Sequence diagram -- public seat lookup through the REST API}
\end{figure}

\subsection{Seating Generation from the Admin Page}
\begin{figure}[H]\centering
\resizebox{\linewidth}{!}{\begin{tikzpicture}[font=\scriptsize,
  head/.style={draw=egnavy,fill=eglight,rounded corners=2pt,minimum width=2.1cm,minimum height=0.65cm,align=center,font=\scriptsize\bfseries},
  msg/.style={-{Stealth[length=2mm]},thick},
  ret/.style={-{Stealth[length=2mm]},dashed}]
\foreach \x/\n/\t in {0/a/Exam Cell\\(Browser),3.1/c/Allocation\\WebController,6.4/s/SeatingAllocation\\Service,9.7/r/Repositories,12.6/d/MySQL}{
  \node[head] (\n) at (\x,0) {\t};
  \draw[eggrey,dashed] (\n.south) -- ++(0,-7.7);
}
\draw[msg] (0,-1.0) -- node[above]{POST /admin/allocation/seating} (3.1,-1.0);
\draw[msg] (3.1,-1.5) -- node[above]{allocateSeats(date, time)} (6.4,-1.5);
\fill[eggold!25] (6.25,-1.7) rectangle (6.55,-5.9);
\node[font=\tiny,rotate=90,eggrey] at (6.75,-3.6) {@Transactional};
\draw[msg] (6.4,-2.0) -- node[above]{findByExamDateAndStartTime} (9.7,-2.0);
\draw[msg] (6.4,-2.5) -- node[above]{findByBranchAndSemester (per exam)} (9.7,-2.5);
\draw[msg] (6.4,-3.0) -- node[above]{hallRepository.findAll} (9.7,-3.0);
\draw[msg] (6.4,-3.5) -- node[above]{deleteSlot (seats, duties)} (9.7,-3.5);
\draw[msg] (9.7,-3.8) -- node[above]{DELETE} (12.6,-3.8);
\draw[msg] (6.4,-4.3) -- ++(0.6,0) -- ++(0,-0.35) node[right]{greedy placement loop} -- ++(-0.6,0);
\draw[msg] (6.4,-5.35) -- node[above]{saveAll(seats)} (9.7,-5.35);
\draw[msg] (9.7,-5.7) -- node[above]{INSERT batch, COMMIT} (12.6,-5.7);
\draw[ret] (6.4,-6.2) -- node[above]{SeatingResult} (3.1,-6.2);
\draw[ret] (3.1,-6.7) -- node[above]{302 redirect /admin/seating + flash} (0,-6.7);
\draw[msg] (0,-7.2) -- node[above]{GET /admin/seating} (3.1,-7.2);
\end{tikzpicture}
}
\caption{Sequence diagram -- generating a seating plan (Post/Redirect/Get)}
\end{figure}
The web controller follows the \emph{Post/Redirect/Get} pattern: after the POST succeeds, the browser is redirected to the plan page, so refreshing the page never repeats the allocation.

\section{User Interface Design}
\subsection{Page Map}
\begin{figure}[H]\centering
\begin{tikzpicture}[font=\scriptsize,
  pg/.style={draw=egnavy,fill=eglight,rounded corners=2pt,minimum width=2.4cm,minimum height=0.6cm,align=center},
  pub/.style={pg,fill=eggold!20,draw=eggold},
  arr/.style={-{Stealth[length=1.8mm]},egnavy}]
\node[pub] (home) at (0,0) {/ (Home)};
\node[pub] (login) at (3.2,0) {/login};
\node[pub] (look) at (-3.2,0) {/seat-lookup};
\node[pg] (dash) at (3.2,-1.4) {/dashboard\\(role redirect)};
\node[pg] (ad) at (0,-2.8) {/admin/dashboard};
\node[pg] (fd) at (5,-2.8) {/faculty/duties};
\node[pg] (sd) at (8,-2.8) {/student/seat};
\foreach \n/\l/\x in {s/students/-4.8,f/faculty/-2.4,h/halls/0,e/exams/2.4,a/allocation/4.8}{
  \node[pg] (\n) at (\x,-4.3) {/admin/\l};
}
\node[pg] (seat) at (3.6,-5.7) {/admin/seating};
\node[pg] (duty) at (6.6,-5.7) {/admin/duties};
\node[pg,fill=white] (pdf) at (5.1,-7.0) {/admin/reports/* (PDF)};
\draw[arr] (home)--(login); \draw[arr] (home)--(look);
\draw[arr] (login)--(dash); \draw[arr] (dash)--(ad); \draw[arr] (dash)--(fd); \draw[arr] (dash)--(sd);
\foreach \n in {s,f,h,e,a} \draw[arr] (ad.south)--(\n.north);
\draw[arr] (a)--(seat); \draw[arr] (a)--(duty); \draw[arr] (seat)--(pdf); \draw[arr] (duty)--(pdf);
\end{tikzpicture}
\caption{Navigation map (orange pages are public)}
\end{figure}

\subsection{Screen Descriptions}
\begin{tabularx}{\linewidth}{|l|X|}
\hline\headrow\thead{Screen} & \thead{Content}\\\hline
Home & Hero banner with Find My Seat and Login buttons; three cards explaining the roles.\\\hline
Login & Username and password; messages for wrong password, logout and expired session.\\\hline
Admin dashboard & Five counters (students, faculty, halls, seats, exams), upcoming slots, workload table.\\\hline
Master data pages & Form on the left (add / edit with validation messages), searchable table on the right with Edit and Delete.\\\hline
Allocation & One card per slot with buttons: 1. Generate seating, 2. Assign invigilators, View plan, Seating PDF.\\\hline
Seating plan & Each hall as a CSS grid; blackboard bar at the top; one pastel colour per subject; empty seats dashed.\\\hline
Duties & Roster table, date filter, PDF button, workload balance table.\\\hline
Seat lookup / My seat & USN search box; result cards with seat number in large font, hall, block, row, bench and timing.\\\hline
My duties & Faculty details and a table filled by a Fetch call to \texttt{/api/duties/faculty/\{code\}}.\\\hline
\end{tabularx}

\section{Security Design}
\begin{tabularx}{\linewidth}{|L{3.4cm}|X|}
\hline\headrow\thead{Concern} & \thead{Design}\\\hline
Authentication & \texttt{AppUserDetailsService} loads \texttt{AppUser} by username (case-insensitive); form login for browsers, HTTP Basic for \texttt{/api/**}.\\\hline
Password storage & \texttt{BCryptPasswordEncoder} (strength 10, random salt).\\\hline
Authorization & URL rules: \texttt{/admin/**} ADMIN, \texttt{/faculty/**} FACULTY, \texttt{/student/**} STUDENT; API: GET \texttt{/api/seats/**} public, \texttt{/api/duties/faculty/**} ADMIN or FACULTY, all other \texttt{/api/**} ADMIN. Extra check in \texttt{DutyApiController} so that a faculty member can read only his or her own roster.\\\hline
Session management & Session fixation protection (\texttt{migrateSession}), one session per user, 30 minute timeout, JSESSIONID deleted on logout.\\\hline
CSRF & Enabled for all browser forms (Thymeleaf adds the token automatically); disabled only for the stateless-style API chain.\\\hline
Output encoding & Thymeleaf \texttt{th:text} escapes HTML; JavaScript builds result cards with an \texttt{escapeHtml} helper or \texttt{textContent}.\\\hline
SQL injection & Only JPA derived queries and parameterised JPQL are used; no string-built SQL.\\\hline
\end{tabularx}

\section{Deployment View}
\begin{figure}[H]\centering
\begin{tikzpicture}[font=\small,
  node3d/.style={draw=egnavy,fill=eglight,thick,minimum width=4.4cm,minimum height=2.4cm,align=center},
  arr/.style={{Stealth[length=2mm]}-{Stealth[length=2mm]},thick,egnavy}]
\node[node3d,fill=white] (c) at (0,0) {\textbf{Client device}\\Browser / Postman};
\node[node3d] (s) at (6,0) {\textbf{Application server}\\JDK 17, examgrid-1.0.0.jar\\Embedded Tomcat :8080};
\node[node3d,fill=egnavy,text=white] (d) at (12,0) {\textbf{Database server}\\MySQL 8 :3306\\examgrid\_db};
\draw[arr] (c)--node[above]{HTTP}(s);
\draw[arr] (s)--node[above]{JDBC}(d);
\end{tikzpicture}
\caption{Deployment diagram (all three nodes may run on one laptop for the demo)}
\end{figure}

\section{Design Decisions Summary}
\begin{tabularx}{\linewidth}{|l|X|X|}
\hline\headrow\thead{No.} & \thead{Decision} & \thead{Reason}\\\hline
D1 & Thymeleaf server-side rendering instead of a SPA framework & Course syllabus; simpler to explain; works without a separate frontend build.\\\hline
D2 & Greedy rule-based allocation instead of backtracking / ILP solver & Runs in linear time, is easy to explain and gives zero conflicts whenever space allows (see Document 03).\\\hline
D3 & DTO records for API responses & Avoids exposing lazy JPA proxies and passwords; stable JSON contract.\\\hline
D4 & Bootstrap through WebJars & Pages work offline in the lab without a CDN.\\\hline
D5 & H2 only in test scope & Tests run anywhere without MySQL; production stays on MySQL.\\\hline
D6 & No Lombok & Plain getters and setters are clearer during a viva.\\\hline
\end{tabularx}

\end{document}
